\chapter{Climate, ecological, planetary failure}
\label{vol10:ch15}

\epigraph{The accident is the entire industrial era; the report is being written by the climate itself.}{}

\chapmeta{Half-life: durable. Archetypes: failure, ethics, scaling. Exercise target: 20. Project track: hybrid.}

\noindent
This chapter places engineering inside the planetary boundaries within which it operates. The failures the earlier chapters have analysed (a stuck valve; a cracked stay; a misclassified bug; a deferred inspection) are local in space and time; the failures this chapter analyses (climate change; stratospheric ozone depletion; lead in petrol; PFAS contamination; ocean acidification; biosphere extinction) are global, multi-decade, and partially irreversible. The engineering profession is a co-author of these failures and a co-author of the responses to them. The chapter is about the discipline of accounting for engineering's planetary footprint, the failure modes that produce planetary-scale damage, and the cases (the ozone success; the lead success; the climate failure in progress; the PFAS failure unfolding) that bound the engineering profession's actual track record.

\section{Engineering as ecological actor}\pathtag{core}

Engineering is a planetary-scale activity. The cumulative mass of human-engineered materials (the "anthropogenic mass": concrete, steel, asphalt, plastics, brick, aluminium, glass) reached approximately the mass of all living biomass on Earth around 2020 by some estimates (Elhacham et al., 2020); the cumulative anthropogenic emissions of greenhouse gases have raised atmospheric CO$_2$ concentrations from approximately 280 ppm at the start of the industrial era to approximately 420 ppm in 2024 \cite{std:ipcc-ar6-wg1-2021}; the engineered footprint of land use (agriculture, urbanisation, infrastructure) covers a substantial fraction of the terrestrial surface; the engineered footprint of fisheries, shipping, and seabed extraction reaches a substantial fraction of the ocean surface and continental shelves.

The framing this chapter takes is that engineering is responsible for the planetary footprint it produces. The framing is not universally accepted within the profession; the alternative framing (engineering is the technical implementation of choices made by clients, regulators, and consumers; the planetary footprint is the result of those choices, not the engineering) has some validity but is incomplete. Engineers make design decisions that shape footprints (materials selection; process efficiency; lifecycle considerations; end-of-life disposition; substitution of harmful with less-harmful technologies). The collective record of these decisions is the engineering profession's planetary footprint; the profession's responsibility is for the part of the footprint that engineering judgment determined.

The disciplines that have grown around this framing include:

\textbf{Life-cycle assessment (LCA).} The methodology for quantifying the environmental impact of a product or process across its life cycle, from raw-material extraction through manufacturing, use, and end-of-life. The ISO 14040 series codifies the methodology; the academic and industry communities have developed databases (ecoinvent, GaBi, USLCI) and software tools (SimaPro, GaBi, OpenLCA) that make LCA accessible. The methodology is the principal engineering tool for quantitative environmental footprint analysis.

\textbf{Industrial ecology.} The discipline that treats industrial systems as analogues of ecological systems: material flows, energy flows, byproducts and wastes treated as inputs to other processes. Yale's Center for Industrial Ecology and the ETH Zürich's industrial ecology program have been the institutional leaders; the discipline has produced extensive analyses of national and global material flows, urban metabolism, and the conditions for closed-loop industrial systems.

\textbf{Green engineering and chemistry.} The design disciplines that incorporate environmental criteria into engineering choices from the beginning rather than as end-of-pipe afterthoughts. The "12 Principles of Green Chemistry" (Anastas and Warner, 1998) and the "12 Principles of Green Engineering" (Anastas and Zimmerman, 2003) are the foundational normative frameworks; their adoption by the chemical and broader engineering industries has been substantial but uneven.

\textbf{Planetary boundaries framework.} The framework (Rockström et al., 2009; updated 2015 and 2023) that identifies nine planetary biophysical thresholds within which human activity can operate sustainably and beyond which the planet's systems become destabilised. The boundaries include climate change, biosphere integrity, biogeochemical flows (nitrogen and phosphorus cycles), land-system change, freshwater use, ocean acidification, atmospheric aerosol loading, stratospheric ozone depletion, and novel entities (chemical pollution). The 2023 assessment identifies six of the nine boundaries as exceeded; the framework is the structured scientific articulation of the claim that planetary-scale failure is in progress.

The engineering profession's relationship to these disciplines is at present partial. LCA is widely used in the chemical, automotive, packaging, and construction industries, with substantial integration into product-design decisions. Industrial ecology is partly absorbed into the engineering curriculum but is not consistently part of the engineer's working toolkit. Green engineering principles are taught in environmental engineering programs and are adopted unevenly in the broader engineering disciplines. The planetary boundaries framework is taught more often in environmental science than in engineering; the connection between specific engineering decisions and the boundaries' state is still being developed.

\begin{principle}[Engineering is accountable for its footprint]
The engineering profession is co-responsible for the planetary footprint of the systems it designs and operates. The framing that locates responsibility elsewhere (in client demand, regulatory permissiveness, consumer preferences) is partial. The engineer's design decisions shape the footprint at the margin and over the long term; the profession's collective record of those decisions constitutes its planetary track record.
\end{principle}

\section{Pollution as engineered byproduct}\pathtag{core}

Pollution is the byproduct that the design did not allocate to a product or to a planned waste stream and that consequently arrives in the environment. The categories of pollution differ in their environmental fate, their human-health consequences, and their reversibility:

\textbf{Air pollutants with short atmospheric residence times.} Sulphur dioxide (residence time days to weeks); nitrogen oxides (similar); particulate matter (days). These pollutants produce regional and local effects (acid deposition; smog; respiratory health impacts) but accumulate only over short times; reductions in emissions produce reductions in atmospheric concentrations within months to years. The success stories of air-pollution control (the post-1970 US Clean Air Act; the Helsinki Protocol on sulphur; the post-1990 acid-rain trading program) demonstrate that short-residence pollutants can be managed effectively through engineering and regulatory action.

\textbf{Air pollutants with long atmospheric residence times.} Carbon dioxide (effective residence time on the order of hundreds to thousands of years for a substantial fraction of any emission); methane (decade); nitrous oxide (century); CFCs (decades to centuries); perfluorocarbons (millennia). These pollutants accumulate over the emission period; reductions in emissions reduce the accumulation rate but do not reduce the accumulated stock (which decays only on the pollutant's natural removal timescale). The management problem is fundamentally different: the engineering choice today affects the planetary state for centuries or millennia.

\textbf{Water pollutants.} Nutrient loading (nitrogen and phosphorus from agricultural runoff and wastewater); heavy metals; synthetic organic compounds; pharmaceutical residues; microplastics; "forever chemicals" (PFAS). The water-quality consequences range from acute (toxic algal blooms; drinking-water contamination) to chronic (endocrine disruption; bioaccumulation in food chains).

\textbf{Soil pollutants.} Heavy metals from industrial activities; persistent organic pollutants (PCBs, dioxins, organochlorine pesticides); residues from mining and extraction; spilled hydrocarbons; PFAS and other forever chemicals. The remediation timescales range from years (bioremediation of degradable compounds) to indefinite (heavy metals; PFAS).

\textbf{Biological pollutants.} Invasive species introduced through trade and transport; cultivars released into wild populations; antimicrobial resistance accumulated in clinical and agricultural use. These pollutants produce cascading effects through ecological networks; reversibility is limited.

The engineering response is the discipline of pollution prevention at design rather than pollution control after release. The hierarchy (preferred to less preferred) is: prevention (the pollution is not produced); minimisation (the pollution is reduced at source); recycling (the pollution is captured and reused); treatment (the pollution is captured and transformed to less harmful form); disposal (the pollution is contained in long-term storage); release (the pollution is emitted to the environment, accepting the consequences). The hierarchy reflects both the lifecycle cost (prevention is usually cheapest over the full lifecycle) and the environmental risk (release is the highest-risk option). The engineering practice of locating pollution-prevention decisions early in design rather than late is the substantive discipline; the alternative (end-of-pipe treatment) is the dominant historical practice and is being progressively replaced.

\section{Climate change engineering-first}\pathtag{core}

Climate change is the engineering profession's largest contemporary failure-and-response. The accumulated CO$_2$ in the atmosphere is the result of approximately two centuries of fossil-fuel combustion; the engineering profession designed and operates the systems that produced the emissions; the engineering profession is principally responsible for designing and deploying the systems that will reduce future emissions and (if the cumulative response is to be adequate) remove past emissions. The chapter's treatment is engineering-first: the climate science is well-developed and is treated as input; the discipline of the chapter is the engineering implications.

The 2021 IPCC Sixth Assessment Report (Working Group I physical-science basis) provides the current authoritative assessment \cite{std:ipcc-ar6-wg1-2021}. The report's central findings: human activities have unequivocally caused approximately $1.1\,\degC$ of global mean surface temperature increase relative to 1850-1900 (the "pre-industrial" baseline); cumulative anthropogenic CO$_2$ emissions through 2019 were approximately 2390 Gt CO$_2$ (give or take ~240 Gt); to limit warming to $1.5\,\degC$ with 50\% probability, cumulative future emissions must be limited to approximately 500 Gt CO$_2$ from 2020 onward (the "remaining carbon budget"); current annual emissions are approximately 40 Gt CO$_2$, so the budget will be exhausted in approximately 12 years at constant emissions; the actual emission trajectory through 2024 is approximately flat to slowly declining, so the actual exhaustion date is some years beyond the constant-emission projection but not dramatically so.

The engineering implications: the global energy system, the global transportation system, the global industrial system, the global building stock, and the global agriculture system must undergo substantial decarbonisation within decades, on a scale and at a speed that has no precedent in the engineering profession's history. The technical scope of the decarbonisation:

\textbf{Electricity.} Approximately 28\% of current global energy use; substantially decarbonisable through wind, solar, nuclear, hydro, and grid energy storage. The technologies are mature; the deployment rate is the rate-limiting factor; the deployment requires substantial investment in transmission, distribution, and storage infrastructure.

\textbf{Industrial heat.} Approximately 25\% of current global energy use; partly decarbonisable through electrification (for lower-temperature heat), partly through hydrogen (for higher-temperature heat), partly through process changes (cement and steel are the largest industrial-emission categories and require process redesign). The technologies are at varying maturity; cement and steel are challenging.

\textbf{Transportation.} Approximately 25\% of current global energy use; substantially decarbonisable through electrification (passenger vehicles; light-duty trucks; rail), with aviation, shipping, and heavy-duty trucking requiring different solutions (electrofuels; hydrogen; biofuels; modal shift).

\textbf{Buildings.} Approximately 20\% of current global energy use; partly through heat pumps, partly through insulation and envelope improvements, partly through changes in operation and use. The existing building stock turnover is slow; substantial retrofit investment is required.

\textbf{Agriculture.} Methane from livestock; nitrous oxide from fertilisers; land-use change; food-system efficiency. The technical and behavioural changes required are substantial; the engineering content includes precision agriculture, alternative protein production, fertiliser-efficiency improvements, and food-supply-chain optimisation.

\textbf{Carbon dioxide removal.} If cumulative emissions exceed the budget for the target temperature, removal of CO$_2$ from the atmosphere becomes necessary to return to the target. Removal technologies (direct air capture; enhanced rock weathering; ocean alkalinity enhancement; biomass with carbon capture and storage; soil-carbon accumulation; afforestation) are at varying maturity and have varying scalability and environmental side-effects.

The engineering profession's role in the decarbonisation is principal. The technical solutions exist or are being developed; the deployment requires engineering work across virtually every engineering discipline; the integration of solutions into existing systems is the substance of the work over coming decades.

The risks of failure (in the sense that the decarbonisation does not occur on the required timescale) include the climate impacts the IPCC identifies (sea-level rise; ecosystem disruption; agricultural disruption; extreme-weather events; the displacement of populations from uninhabitable regions; the loss of biodiversity; the cascade effects through interconnected systems). The risks rise nonlinearly with temperature; the cumulative warming exceeds the boundary at which the planet's systems become substantially destabilised somewhere between $1.5$ and $3\,\degC$, with substantial uncertainty.

\begin{archetype}[scaling]
Climate change is the scaling archetype's planetary expression. The engineering activity at the scale of individual projects is locally manageable; the cumulative footprint at planetary scale is the failure. The disconnect between the local optimisation problem each engineer solves and the global integrated effect is the conditions under which the failure has accumulated.
\end{archetype}

\section{Tipping points and irreversibility}\pathtag{standard}

The climate and ecological systems are not linear: substantial regions of state space are bistable, hysteretic, or characterised by tipping points beyond which the system reorganises into a different stable state from which return is difficult or impossible on relevant timescales.

The major climate tipping points identified in the current literature:

\textbf{Greenland ice sheet collapse.} The ice sheet's mass loss is currently accelerating; sustained warming above some threshold (estimates centre near $1.5$--$2\,\degC$) appears to commit the ice sheet to long-term complete loss, contributing approximately 7 m of eventual sea-level rise over multiple millennia.

\textbf{West Antarctic Ice Sheet collapse.} Marine-terminating glaciers in the Amundsen Sea sector appear already to have crossed the threshold of irreversible retreat; the eventual sea-level contribution is approximately 3 m, over centuries.

\textbf{Atlantic Meridional Overturning Circulation (AMOC) slowdown or collapse.} The AMOC drives substantial heat transport from low to high latitudes in the Atlantic; its weakening would produce substantial regional climate changes (cooler Western Europe; shifted tropical rainfall patterns); the threshold is uncertain but recent observational evidence suggests substantial weakening is already occurring.

\textbf{Amazon rainforest dieback.} Sustained warming combined with land-use change appears to threaten the rainforest's stability; the loss would shift the region from a carbon sink to a carbon source and would have substantial regional climate consequences.

\textbf{Permafrost carbon release.} Thawing permafrost releases CO$_2$ and methane from previously-frozen organic matter; the release is a positive feedback on the warming that thaws it; the threshold for substantial accelerated release is uncertain.

\textbf{Coral reef collapse.} Warming and ocean acidification together are pushing coral reefs past their thermal and chemical tolerance; the dominant tropical-reef ecosystems appear unlikely to persist in their current form even under aggressive warming-limitation scenarios.

The engineering implications of tipping-point behaviour: traditional engineering risk analysis assumes that more cautious operation reduces risk smoothly; tipping-point behaviour means that operation that approaches a tipping point can be safe up until the threshold is crossed, after which the consequences are large and irreversible. The risk profile is more demanding: the precautionary principle (avoid approaches that could cross tipping points even if the costs of avoidance are substantial) becomes the engineering basis rather than an external ethical add-on.

The irreversibility dimension is particularly relevant. A water-pollution problem can be remediated over years; an emissions excursion can be reversed over decades; a tipping-point crossing cannot be reversed on engineering timescales. The engineering disciplines that have developed to operate within hard limits (radiation-safety regulation in nuclear; emissions caps in air-quality regulation; quotas in fisheries management) are partial models for what climate and ecological engineering require at planetary scale.

\section{Engineering responses to ecological failure}\pathtag{standard}

The engineering responses to planetary-scale ecological failure operate at several levels:

\textbf{Substitution of harmful technologies.} Replace the technology that produces the harm with one that produces less. The decarbonisation of electricity, the replacement of leaded petrol with unleaded, the replacement of CFCs with HFCs (and subsequently with HFOs and natural refrigerants) are the canonical successful examples. The substitution requires alternatives that are technically feasible, economically viable, and deployable at scale; the engineering work is the development and deployment of the alternatives.

\textbf{Efficiency improvement.} Reduce the resource intensity per unit of service delivered. The improvements in energy efficiency of buildings, vehicles, appliances, and industrial processes over recent decades have reduced the resource intensity substantially; the further improvements available are bounded by physics but extend substantial distances at current technology levels.

\textbf{Recycling and circular-economy redesign.} Move from linear material flows (extraction, use, disposal) to circular flows (use, recovery, reuse). The circular-economy framing has been adopted in some industries (aluminium, steel, paper, glass) and is being extended into others (electronics, plastics, batteries, construction materials). The engineering work is the design of products for disassembly and recovery, the design of recovery processes, and the integration of recovered materials into new production.

\textbf{Pollution control and treatment.} Capture the harmful byproducts and treat them to less-harmful forms or contain them in long-term storage. The technologies range from mature (flue-gas desulphurisation; catalytic converters; municipal wastewater treatment) to emerging (carbon capture and storage; PFAS destruction). Pollution control is the historical workhorse of environmental engineering and remains important even as the discipline shifts toward prevention.

\textbf{Carbon dioxide removal.} For the climate problem specifically, the cumulative emissions already in the atmosphere will produce substantial warming even with zero further emissions; if the target is to stabilise warming at a level below the current trajectory, CO$_2$ removal becomes part of the engineering response. The technologies are at varying maturity; their deployment at the scale required (gigatons per year) raises substantial questions about energy, land, and water footprints.

\textbf{Geoengineering.} Deliberate large-scale intervention in Earth-system processes (stratospheric aerosol injection; marine cloud brightening; cirrus thinning; ocean fertilisation) to reduce warming. The interventions raise governance, ethical, and unintended-consequences questions that make their deployment contested; the engineering work on geoengineering is at present primarily research-stage with limited deployment.

\textbf{Adaptation.} Where mitigation does not prevent the consequences, adapt the engineered environment to operate under the changed conditions: coastal protection against sea-level rise; air conditioning capacity for increased heat; drought-tolerant infrastructure in regions becoming arid; flood-tolerant infrastructure in regions becoming wet. Adaptation is the engineering work of operating under failed mitigation; it is not a substitute for mitigation but a necessary complement.

The engineering profession's contribution to the responses is the work of decades; the cumulative effect will determine whether the planetary-scale failure becomes catastrophic or merely substantial.

\section{Case: the ozone-layer near-miss (and the success), CFC industry, lead in petrol, PFAS}\pathtag{core}

\subsection*{The ozone-layer near-miss and the Montreal Protocol success}

The stratospheric ozone layer protects the Earth's surface from harmful ultraviolet radiation. In the early 1970s, Mario Molina and F. Sherwood Rowland published the theoretical basis for the destruction of stratospheric ozone by chlorofluorocarbon (CFC) refrigerants and propellants \cite{paper:molina-rowland-1974}: CFCs, chemically inert in the troposphere and with atmospheric residence times of decades to centuries, reach the stratosphere where ultraviolet radiation breaks them apart, releasing chlorine atoms that catalyse the destruction of ozone molecules in a chain reaction. The theoretical prediction was that sustained CFC release would substantially reduce stratospheric ozone over decades to centuries.

The theoretical prediction was substantially confirmed by the discovery of the Antarctic ozone hole by Joseph Farman, Brian Gardiner, and Jonathan Shanklin in 1985 \cite{paper:farman-1985}. The observational data from the British Antarctic Survey's Halley research station showed substantial seasonal depletion of stratospheric ozone over Antarctica during austral spring, with the depletion having begun in the mid-1970s and increased through the 1980s. The observational discovery accelerated the international response.

The Montreal Protocol on Substances that Deplete the Ozone Layer was adopted in 1987 and entered into force in 1989 \cite{law:montreal-protocol-1987}. The Protocol's structure was a phased schedule for the reduction and eventual elimination of CFC production, applied to all parties (with longer compliance schedules for developing countries). Subsequent amendments (London 1990; Copenhagen 1992; Kigali 2016) tightened the schedule and added additional ozone-depleting substances and (in Kigali) the high-global-warming-potential HFCs that had replaced CFCs.

The engineering response was the development and deployment of CFC alternatives. HCFCs (hydrochlorofluorocarbons) were the first-generation substitutes; they have lower ozone-depleting potential than CFCs but are not zero. HFCs (hydrofluorocarbons) are second-generation substitutes; they have zero ozone-depleting potential but high global-warming potential. HFOs (hydrofluoroolefins) and natural refrigerants (ammonia, CO$_2$, hydrocarbons) are third-generation substitutes; they have zero or very low ozone-depleting potential and low to moderate global-warming potential. The substitution sequence reflects accumulating learning about the full environmental consequences of the alternatives.

The outcome is among the engineering profession's most direct planetary-scale successes. Atmospheric CFC concentrations peaked in the mid-1990s and are now in slow decline; the Antarctic ozone hole has stabilised and shows evidence of recovery; the projection is that the ozone layer will return to pre-1980 levels by approximately 2050-2070. The success has been achieved through international institutional cooperation (the Protocol and its amendments), engineering substitution at industrial scale, and the technical and economic feasibility of the alternatives. The pattern (theoretical prediction, observational confirmation, international response, engineering substitution, planetary-scale recovery) is the engineering profession's working model for what planetary-scale environmental response can look like.

The "near-miss" framing applies to what would have happened in the absence of the response. The CFC release trajectory in the early 1980s, projected forward without the Montreal Protocol, would have produced severe stratospheric ozone depletion globally by the mid-21st century, with associated increases in UV radiation, skin cancers, ecosystem disruption, and broader climate effects. The response prevented the projected damage; the lesson is that planetary-scale failure can be averted when the engineering profession, the political institutions, and the affected industries align on the response.

\subsection*{Lead in petrol}

Tetraethyl lead (TEL) was added to gasoline beginning in the 1920s as an octane booster; its use spread globally; cumulative atmospheric lead emissions from leaded petrol over the 20th century contaminated soils, water, and (most consequentially) human populations worldwide. The toxic effects of lead on the human nervous system, particularly in children, are substantial and were extensively documented in the public-health literature over decades \cite{paper:needleman-lead-1990}. The phase-out of leaded petrol began in the US in the 1970s and 1980s under Clean Air Act regulations and was completed for road fuels worldwide in the early 2020s (Algeria, the last country, transitioned in 2021).

The engineering response was the development of alternative octane boosters (MMT, MTBE, ethanol, alkylates) and the redesign of engines to operate without the lead-tolerant valve seats. The substitution was complete for road fuels; aviation fuel (avgas, used in piston-engine aircraft) still contains lead and is the subject of ongoing engineering work on substitutes.

The lead-in-petrol case is structurally similar to the ozone case in some respects (a substance recognised as harmful; a substitution pathway developed; a phased global elimination achieved) but the timescale was longer (decades rather than years for substantial action), the institutional response was less centralised (national regulation rather than international protocol), and the engineering substitution faced additional barriers (the engine-redesign requirement). The case demonstrates that the engineering profession can complete a difficult substitution at planetary scale, even when the substitution takes longer than ideal; the cumulative human-health benefit is substantial.

The case also demonstrates the role of public-health science in driving engineering response. The Needleman et al. work on the long-term effects of low-level lead exposure on children's neurodevelopment was contested by the industry at the time; the resolution of the controversy in favour of the public-health evidence shaped the regulatory case for the phase-out. The lesson for engineering is partly about the connection between engineering substitution and the scientific evidence base that supports the regulatory case.

\subsection*{PFAS: the failure in progress}

Per- and polyfluoroalkyl substances (PFAS) are a family of approximately 10,000 to 14,000 synthetic compounds characterised by very strong carbon-fluorine bonds that confer thermal and chemical stability; the same property that makes them useful in industrial applications (non-stick coatings; firefighting foams; water-repellent textiles; semiconductor manufacturing; food packaging) makes them resistant to environmental degradation. Some PFAS compounds have atmospheric and aquatic residence times measured in centuries to millennia; they accumulate in soils, water, and biological tissue across the planet; the human-health consequences (developmental toxicity; immune-system effects; thyroid effects; some cancers) are substantial and continue to be characterised.

The engineering and regulatory response is in progress and is incomplete. The two best-characterised PFAS compounds (PFOA, used in Teflon manufacture; PFOS, used in firefighting foam and stain repellents) have been phased out in major markets through the Stockholm Convention and national regulations; the substitution has often been to other PFAS compounds (shorter-chain perfluoroalkyl acids; fluoropolymer alternatives) whose long-term consequences are not yet fully characterised. The detection of PFAS contamination in drinking-water supplies, in surface waters, in soils, in biota, and in human blood serum (essentially every measured population has detectable levels) is now routine; the engineering and remediation challenge of removing PFAS from contaminated water and from the broader environment is substantial.

The PFAS case is the current canonical example of a planetary-scale engineering failure unfolding in real time. The compounds were developed and deployed at industrial scale before their environmental persistence and toxicity were characterised; the substitution to alternatives has often been incomplete (shorter-chain PFAS may be substantially as persistent as their predecessors); the remediation of contaminated environments is at very early stages; the engineering and regulatory response is uneven across jurisdictions. The case illustrates the failure mode in which the engineering profession deployed a class of substances at planetary scale without adequate prior characterisation of their environmental fate, and is now in the position of attempting to manage the consequences after the deployment has occurred.

\begin{principle}[The engineering response to ecological failure can succeed]
The ozone and lead cases demonstrate that the engineering profession can respond effectively to planetary-scale environmental failures when the conditions align: scientific characterisation of the harm; institutional frameworks for international or coordinated response; technical and economic feasibility of substitutes; sustained engineering work over the decades required. The conditions are not the default; they have to be constructed. The cases in which they were constructed are the engineering profession's most consequential planetary contributions; the cases in which they were not (PFAS in progress; climate the largest) are the failures the next decades of engineering work must address.
\end{principle}

\begin{project}[Environmental footprint assessment of one product]
Select one product the reader uses (a piece of clothing; a phone; a vehicle; a building material; a packaged food). Conduct a streamlined life-cycle assessment using primary data where available and best-available secondary data otherwise. The assessment covers: raw-material extraction (what materials, from where, with what footprint); manufacturing (energy and material inputs; emissions; waste streams); transportation (from manufacture to use); use phase (energy and material consumption during use; emissions); end of life (recovery, recycling, disposal). Report the dominant impact category (climate; water; toxicity; land use; resource depletion) and the dominant life-cycle stage. Identify two engineering changes that would substantially reduce the footprint and assess the feasibility of each. Deliverable: a four-page document with the assessment, the impact summary, and the recommendations.

\textbf{Hazard.} The exercise often reveals that the dominant impact category and stage are not what the reader assumed (use-phase energy for some products; manufacturing for others; supply-chain transport for others). The point is to develop the engineering habit of footprint analysis as an input to design decisions; the analysis is not advocacy but engineering work.
\end{project}

\section{Exercises}

\subsection*{Calculation}\pathtag{core}

\begin{exercise}
Atmospheric CO$_2$ concentration was approximately 280 ppm in 1850 and approximately 420 ppm in 2024 \cite{std:ipcc-ar6-wg1-2021}. Compute the percentage increase. Estimate the cumulative anthropogenic emissions implied by the increase (using a conversion factor of approximately 7.8 Gt CO$_2$ per ppm increase, accounting for ocean and biosphere uptake of approximately half of cumulative emissions).
\end{exercise}

\begin{exercise}
The remaining carbon budget for $1.5\,\degC$ with 50\% probability is approximately 500 Gt CO$_2$ from 2020 onward. Annual emissions in 2024 are approximately 40 Gt CO$_2$. Compute the years remaining at constant emissions. If emissions decline linearly to zero over $T$ years, compute the cumulative emissions and solve for the $T$ that exhausts the budget.
\end{exercise}

\begin{exercise}
A typical European household consumes approximately 3500 kWh of electricity per year. The carbon intensity of the EU electricity mix in 2024 is approximately 250 g CO$_2$/kWh. Compute the household's annual emissions from electricity. Discuss the dependence on the mix's carbon intensity, and identify the conditions under which decarbonisation of electricity supply reduces household emissions versus the conditions under which behavioural changes do.
\end{exercise}

\begin{exercise}
The atmospheric lifetime of CFC-12 is approximately 100 years. Compute the fraction remaining after 50 years and after 500 years of zero emissions. Discuss the implication for the timescale of stratospheric ozone recovery following the Montreal Protocol \cite{law:montreal-protocol-1987}.
\end{exercise}

\subsection*{Derivation}\pathtag{standard}

\begin{exercise}
Derive the carbon-cycle equilibrium expression: given anthropogenic emissions $E(t)$, derive the atmospheric CO$_2$ concentration $C(t)$ as a function of $E$ and the carbon-removal rates (ocean uptake; terrestrial biosphere uptake). Identify the conditions under which atmospheric CO$_2$ stabilises versus continues to rise.
\end{exercise}

\begin{exercise}
Derive a simple tipping-point model: a system with state $S$ has potential $V(S) = -aS^2 + bS^4 + cS$ where $c$ is an external forcing. Identify the values of $c$ at which the potential has one stable minimum versus two. Identify the hysteresis loop the system traces as $c$ varies cyclically.
\end{exercise}

\begin{exercise}
For a pollution-prevention hierarchy with prevention cost $p$, treatment cost $t$, and damage cost (if released) $d$, derive the conditions under which prevention is preferred to treatment and under which both are preferred to release. Identify the role of discounting in the analysis.
\end{exercise}

\subsection*{Estimation}\pathtag{standard}

\begin{exercise}
Estimate the energy required to remove 1 Gt of CO$_2$ from the atmosphere by direct air capture using current technology. Compare to current global annual primary energy consumption. Discuss the implications for the scalability of direct air capture as a climate response.
\end{exercise}

\begin{exercise}
Estimate the cumulative lead deposition from petrol combustion in a typical mid-20th-century US urban environment over the 1950-1990 period. Discuss the implications for current soil-lead burden and the engineering remediation challenge.
\end{exercise}

\begin{exercise}
Estimate the cost of replacing all PFAS-contaminated drinking-water supplies in a typical US municipality with $10^5$ residents, assuming the contamination requires activated-carbon or ion-exchange treatment at the treatment plant.
\end{exercise}

\subsection*{Simulation}\pathtag{standard}

\begin{exercise}
Implement a simple climate-economy model: atmospheric CO$_2$ concentration evolves with emissions and natural sinks; temperature responds to concentration with a time constant; emissions evolve with GDP and carbon intensity; carbon intensity declines with policy and technology. Sweep policy strength and identify the policy regime that limits warming to $2\,\degC$.
\end{exercise}

\begin{exercise}
Simulate the Antarctic ozone hole's recovery under the Montreal Protocol's actual phase-out schedule and a counterfactual no-Protocol scenario \cite{law:montreal-protocol-1987,paper:farman-1985}. Identify the magnitude of stratospheric ozone the Protocol preserved.
\end{exercise}

\subsection*{Design}\pathtag{standard}

\begin{exercise}
Design a streamlined life-cycle assessment template for use in early-stage engineering design decisions. Specify the impact categories, the data sources, the boundary conditions, and the format for communicating results to non-LCA-specialist designers.
\end{exercise}

\begin{exercise}
Design an engineering substitution roadmap for a hypothetical industrial process currently using a hazardous chemical (e.g., a chlorinated solvent in metal degreasing). Specify the alternatives, the technical evaluation criteria, the deployment schedule, and the verification regime.
\end{exercise}

\begin{exercise}
Design an adaptation strategy for a coastal city facing $1\,\text{m}$ of sea-level rise over the next century. Specify the engineering interventions (sea walls; managed retreat; elevated infrastructure; stormwater management), the priorities, the timing, and the funding mechanism.
\end{exercise}

\subsection*{Judgement}\pathtag{mastery}

\begin{exercise}
Argue: the engineering profession's principal responsibility regarding climate change is to develop and deploy decarbonisation technologies, with the political and economic decisions about deployment being outside the engineering scope. Argue the opposite. Identify what the profession's codes of ethics actually support.
\end{exercise}

\begin{exercise}
Discuss the structural conditions that produced the Montreal Protocol success \cite{law:montreal-protocol-1987} versus the conditions that have produced the slower climate response. Identify the differences that are intrinsic to the problems (scale, complexity, economic stakes) versus those that are contingent on institutional choices.
\end{exercise}

\begin{exercise}
Argue: the PFAS case demonstrates that the engineering profession should not deploy any new class of persistent synthetic chemicals at industrial scale without prior characterisation of long-term environmental fate. Identify the practical constraints on the position and discuss what regulatory and professional structures would implement it.
\end{exercise}

\begin{exercise}
Geoengineering proposals (stratospheric aerosol injection; marine cloud brightening) are technically feasible but raise substantial governance, ethical, and unintended-consequences questions. Discuss the conditions under which the engineering profession should research and prepare such technologies versus the conditions under which research itself is irresponsible.
\end{exercise}

\begin{exercise}
The cumulative engineering work required for adequate climate response over coming decades exceeds the scale of any prior engineering programme. Discuss whether the engineering profession's current institutional structure (national professional bodies; academic curricula; industry employment) is adequate to the scale, and identify the changes that would substantially increase the profession's capacity.
\end{exercise}
